
\documentclass[12pt]{article}
\usepackage[utf8]{inputenc}
\usepackage{graphicx}
\usepackage{float}
\usepackage{listings}
\usepackage{xcolor}
\usepackage{geometry}
\geometry{margin=1in}

\title{Aplicación en Qt5: Publicador y Suscriptor de URL de Video Embebido}
\author{Curso ELO329 - Diseño y Programación Orientados a Objetos}
\date{}

\definecolor{codebg}{rgb}{0.95,0.95,0.95}
\lstset{
  backgroundcolor=\color{codebg},
  basicstyle=\ttfamily\small,
  breaklines=true,
  frame=single,
  postbreak=\mbox{\textcolor{red}{$\hookrightarrow$}\space},
  language=C++,
  showstringspaces=false
}

\begin{document}

\maketitle

\section*{Descripción General}

Se presenta una aplicación Qt5 basada en el patrón \textbf{Publisher/Subscriber} que permite publicar y visualizar URLs de videos. A diferencia de la versión básica que abre los videos en un navegador externo, esta implementación embebe el video directamente en la aplicación utilizando \texttt{QWebEngineView}.

\begin{itemize}
  \item \textbf{Publicador}: El usuario ingresa una URL de video en un campo de texto y la publica mediante un botón.
  \item \textbf{Suscriptor}: Recibe la última URL publicada y al hacer clic en el botón, reproduce el video embebido dentro de una ventana Qt.
\end{itemize}

\section*{Código Fuente Principal}

\subsection*{subscriberwindow.h}
\begin{lstlisting}
class SubscriberWindow : public QMainWindow {
    Q_OBJECT
public:
    explicit SubscriberWindow(QWidget *parent = nullptr);
    ~SubscriberWindow();

public slots:
    void updateUrl(const QString &url);

private slots:
    void on_openVideoButton_clicked();

private:
    Ui::SubscriberWindow *ui;
    QString currentUrl;
    QWebEngineView *videoView;
};
\end{lstlisting}

\subsection*{subscriberwindow.cpp}
\begin{lstlisting}
#include "subscriberwindow.h"
#include "ui_subscriberwindow.h"
#include <QWebEngineView>

SubscriberWindow::SubscriberWindow(QWidget *parent)
    : QMainWindow(parent)
    , ui(new Ui::SubscriberWindow) {
    ui->setupUi(this);
    currentUrl = "";
    ui->openVideoButton->setEnabled(false);
    videoView = nullptr;
}

SubscriberWindow::~SubscriberWindow() {
    delete ui;
}

void SubscriberWindow::updateUrl(const QString &url) {
    currentUrl = url;
    ui->openVideoButton->setText("Ver video");
    ui->openVideoButton->setEnabled(true);
}

void SubscriberWindow::on_openVideoButton_clicked() {
    if (!currentUrl.isEmpty()) {
        if (!videoView) {
            videoView = new QWebEngineView(this);
            videoView->resize(800, 600);
        }
        videoView->load(QUrl(currentUrl));
        videoView->show();
    }
}
\end{lstlisting}

\subsection*{QtVideoPublisher.pro}
\begin{lstlisting}[language=bash]
QT += core gui widgets webenginewidgets

TARGET = QtVideoPublisher
TEMPLATE = app

SOURCES += Sources/main.cpp \
           Sources/publisherwindow.cpp \
           Sources/subscriberwindow.cpp

HEADERS += Headers/publisherwindow.h \
           Headers/subscriberwindow.h

FORMS += Forms/publisherwindow.ui \
         Forms/subscriberwindow.ui
\end{lstlisting}

\noindent
\textbf{Nota:} Se ha agregado \texttt{webenginewidgets} al archivo \texttt{.pro} para habilitar el uso de \texttt{QWebEngineView}.

\end{document}
